
This appendix presents the proofs for the unified framework \autoref{thm:unified_obj_main_text}. We show that neuron addition methods
such as TINY~\citep{verbockhaven_growing_2024}, GradMax~\citep{evci_gradmax_2022}, SENN~\citep{mitchell_self-expanding_2024}, and NeST~\citep{daiNeSTNeuralNetwork2019} can all be understood as special
cases of a common optimization objective, differing only in their choice
of metric and constraints.

\subsection{Preliminaries}

\begin{lemma}\label{lemma:argument_scaling}

For:

\begin{itemize}
\item \(E\) a pre-Hilbert space,
\item \(F\) a non empty set stable by scaling \emph{i.e.}
  (\(\forall x \in F,\forall\lambda \in \mathbb{R},\lambda x \in F\))
\item \(f:F \rightarrow E\) respect some scaling homogeneity
  (\(\exists p \in \mathbb{N}^{\ast},\forall x \in F,\forall\lambda \in \mathbb{R},f(\lambda x) = \lambda^{p}f(x)\))
\item \(y \in E\).
\end{itemize}
If \(\left( \exists x \in F,{\langle f(x),y\rangle} > 0 \right)\) we
have that:
\[\argmax_{d \in F~|~\left\| {f(d)} \right\| \leq 1}{\langle y,f(d)\rangle} = \argmax_{d \in F~|~\left\| {f(d)} \right\| = 1}{\langle y,f(d)\rangle}\]
and:
\[\argmin_{x \in F}\left\| {y - f(x)} \right\| = \argmin_{x \in F~|~\left\| {f(x)} \right\| \neq 0}\left\| {y - f(x)} \right\|\]

\end{lemma}

\begin{proof}

Let
\(I \coloneqq \argmax_{d \in F|\left\| {f(d)} \right\| \leq 1}{\langle y,f(d)\rangle}\)
and
\(R \coloneqq \argmax_{d \in F|\left\| {f(d)} \right\| = 1}{\langle y,f(d)\rangle}\).

\begin{itemize}
\item
  Let us show that \(I \subset R\). Let \(d^{\ast} \in I\); we will show that
  \(\left\| {f\left( d^{\ast} \right)} \right\| = 1\).

  Let \(x \in F\) with \({\langle f(x),y\rangle} > 0\) (exists by
  hypothesis). For \(\varepsilon > 0\) small enough,
  \(\left\| {f(\varepsilon x)} \right\| = \varepsilon^{p}\left\| {f(x)} \right\| \leq 1\),
  so \(\varepsilon x\) satisfies the norm constraint in the definition of \(I\), and
  \({\langle y,f(\varepsilon x)\rangle} = \varepsilon^{p}{\langle y,f(x)\rangle} > 0\).
  By optimality of \(d^{\ast}\) in \(I\),
  \({\langle y,f\left( d^{\ast} \right)\rangle} \geq {\langle y,f(\varepsilon x)\rangle} > 0\),
  hence \(f\left( d^{\ast} \right) \neq 0\). Let
  \(\lambda \coloneqq \frac{1}{\left\| {f\left( d^{\ast} \right)} \right\|^{\frac{1}{p}}} > 0\).
  As \(|\lambda| \geq 1\):
  \[{\langle y,f\left( d^{\ast} \right)\rangle} \leq \lambda^{p}{\langle y,f\left( d^{\ast} \right)\rangle} = {\langle y,f\left( \lambda d^{\ast} \right)\rangle}\]
  by homogeneity of \(f\). By maximality of \(d^{\ast}\) in \(I\), as
  \(\left\| {f\left( \lambda d^{\ast} \right)} \right\| = 1\), this is
  an equality hence \(\lambda = 1\) and
  \(\left\| {f\left( d^{\ast} \right)} \right\| = 1\). Therefore
  \(d^{\ast} \in R\) and \(I \subset R\).
\item
  Let us show that \(R \subset I\). Let \(d^{\ast} \in R\); since
  \(I \subset R\), \(d^{\ast}\) maximizes the score over
  \(\left\{ d \in F~|~\left\| {f(d)} \right\| \leq 1 \right\}\), and as
  \(\left\| {f\left( d^{\ast} \right)} \right\| \leq 1\) we have that
  \(d^{\ast} \in I\). We conclude that \(R \subset I\).
\end{itemize}

For the second equality, with the witness \(x\) above,
\(\left\| {y - f(\varepsilon x)} \right\|^{2} = \left\| y \right\|^{2} - 2\varepsilon^{p}{\langle y,f(x)\rangle} + \varepsilon^{2p}\left\| {f(x)} \right\|^{2} < \left\| y \right\|^{2}\)
for \(\varepsilon > 0\) small, while \(f(x') = 0\) gives
\(\left\| {y - f(x')} \right\|^{2} = \left\| y \right\|^{2}\); so no
minimizer of \(\left\| {y - f(x)} \right\|\) has \(f(x) = 0\).

\end{proof}

\begin{lemma}\label{lemma:sign_symmetrization}

For:

\begin{itemize}
\item
  \(E\) a pre-Hilbert space,
\item
  \(F\) a non-empty set,
\item
  \(f:F \rightarrow E\),
\item
  \(y \in E\),
\item
  \(n:F \rightarrow F\) a sign-flip operator:
  \(\forall d \in F,f\left( n(d) \right) = - f(d)\).
\end{itemize}

Then \(\langle f\left( . \right),y\rangle\) takes a strictly positive
value as soon as it takes a non-zero one:
\[\left( \exists x \in F,{\langle f(x),y\rangle} \neq 0 \right) \Longrightarrow \left( \exists x \in F,{\langle f(x),y\rangle} > 0 \right).\]

\end{lemma}

\begin{proof}

Let \(x \in F\) with \({\langle f(x),y\rangle} \neq 0\). If
\({\langle f(x),y\rangle} > 0\) we are done; otherwise
\({\langle f(x),y\rangle} < 0\) and \(n(x) \in F\) satisfies
\({\langle f\left( n(x) \right),y\rangle} = - {\langle f(x),y\rangle} > 0\).

\end{proof}

\begin{theorem}\label{lemma:argmin_norm_to_argmax_scalarp}

For:

\begin{itemize}
\item
  \(E\) a pre-Hilbert space,
\item
  \(F\) a non empty set stable by scaling \emph{i.e.}
  (\(\forall x \in F,\forall\lambda \in \mathbb{R},\lambda x \in F\))
\item
  \(f:F \rightarrow E\) respect some scaling homogeneity
  (\(\exists p \in \mathbb{N}^{\ast},\forall x \in F,\forall\lambda \in \mathbb{R},f(\lambda x) = \lambda^{p}f(x)\))
\item
  \(n:F \rightarrow F\) a sign-flip operator:
  \(\forall d \in F,f\left( n(d) \right) = - f(d)\)
\item
  \(y \in E\) such that
  \(\exists x \in F,{\langle f(x),y\rangle} \neq 0\) (non-degeneracy).
\end{itemize}

We have that:

\begin{equation}\label{eq:fixed_norm_optimization_equivalence}
\argmin_{x \in F}\left\| {y - f(x)} \right\| = \left\{ {\langle y,f\left( d^{\ast} \right)\rangle}^{\frac{1}{p}}d^{\ast}~|~d^{\ast} \in \argmax_{d \in F~|~\left\| {f(d)} \right\| \leq 1}{\langle y,f(d)\rangle} \right\}
\end{equation}

\end{theorem}

\begin{proof}

By \autoref{lemma:sign_symmetrization} the non-degeneracy hypothesis yields
some \(x \in F\) with \({\langle f(x),y\rangle} > 0\), so by
\autoref{lemma:argument_scaling} we can add the constraint
\(\left\| {f(x)} \right\| \neq 0\) in the left-hand side and replace
\(\leq 1\) by \(= 1\) in the right-hand side. Let's proceed by setting
equality. \[\begin{aligned}
\argmin_{x \in F~|~\left\| {f(x)} \right\| \neq 0}\left\| {y - f(x)} \right\| & = \argmin_{x \in F}\left( \left\| y \right\|^{2} - 2{\langle y,f(x)\rangle}+ \left\| {f(x)} \right\|^{2}  \right) \\
 & = \argmin_{x \in F~|~\left\| {f(x)} \right\| \neq 0}\left( - 2{\langle y,f(x)\rangle} + \left\| {f(x)} \right\|^{2} \right) \\
 & = \argmin_{x \in F~|~\left\| {f(x)} \right\| \neq 0}\left( - 2{\langle y,\left\| {f(x)} \right\|\frac{f(x)}{\left\| {f(x)} \right\|}\rangle} + \left\| {f(x)} \right\|^{2} \right) \\
\text{As}\ f\ \text{is}\ \text{homogenous}: & = \argmin_{x \in F~|~\left\| {f(x)} \right\| \neq 0}\left( - 2\left\| {f(x)} \right\|{\langle y,f\left( \frac{x}{\left\| {f(x)} \right\|^{\frac{1}{p}}} \right)\rangle} + \left\| {f(x)} \right\|^{2} \right) \\
 & = \argmin_{x \in F~|~\left\| {f(x)} \right\| \neq 0}\left( - 2t{\langle y,f(d)\rangle} + t^{2} \right)
\end{aligned}\] with
\(d \coloneqq \frac{x}{\left\| {f(x)} \right\|^{\frac{1}{p}}} \in F\)
and \(t \coloneqq \left\| {f(x)} \right\| \geq 0\). The solution over
\(t\) is achieved at \(t = {\langle y,f(d)\rangle}\). Hence we get:
\[\begin{aligned}
\argmin_{x \in F~|~\left\| {f(x)} \right\| \neq 0}\left\| {y - f(x)} \right\| & = \argmin_{x \in F~|~\left\| {f(x)} \right\| \neq 0}\left( - {\langle y,f(d)\rangle}^{2} \right) \\
 & = \argmax_{x \in F~|~\left\| {f(x)} \right\| \neq 0}{\langle y,f(d)\rangle} \\
 & = \left\{ {\langle y,f\left( d^{\ast} \right)\rangle}^{\frac{1}{p}}d^{\ast}~|~d^{\ast} \in \argmax_{d \in F~|~\left\| {f(d)} \right\| = 1}{\langle y,f(d)\rangle} \right\}
\end{aligned}\]

\end{proof}

\begin{corollary}\label{corollary:argmin_norm_to_argmax_scalarp}

For:

\begin{itemize}
\item
  \(E\) a pre-Hilbert space,
\item
  \(F\) a non empty set stable by scaling \emph{i.e.}
  (\(\forall x \in F,\forall\lambda \in \mathbb{R},\lambda x \in F\))
\item
  \(f:F \rightarrow E\) respect some scaling homogeneity
  (\(\exists p \in \mathbb{N}^{\ast},\forall x \in F,\forall\lambda \in \mathbb{R},f(\lambda x) = \lambda^{p}f(x)\))
\item
  \(n:F \rightarrow F\) a sign-flip operator:
  \(\forall d \in F,f\left( n(d) \right) = - f(d)\)
\item
  \(y \in E\) such that
  \(\left( \exists x \in F,{\langle f(x),y\rangle} \neq 0 \right)\) % or
%  \(\ker f = \left\{ 0 \right\}\).
\end{itemize}

We have that:
\[\argmin_{x \in F}\left\| {y - f(x)} \right\| \propto \argmax_{d \in F|\left\| {f(d)} \right\| \leq 1}{\langle y,f(d)\rangle}\]

\end{corollary}

\begin{proof}
\leavevmode\par\vspace{0.5\baselineskip}

If \(\exists x \in F,{\langle f(x),y\rangle} \neq 0\), the non-degeneracy hypothesis of \autoref{lemma:argmin_norm_to_argmax_scalarp} is met (using the same \(n\)), and its right-hand side is exactly the desired rescaled right-hand \(\argmax\).

% \begin{itemize}
%    \item Otherwise \(\forall x \in F,{\langle f(x),y\rangle} = 0\) and \(\ker f = \left\{ 0 \right\}\). Then \(\left\| {y - f(x)} \right\|^{2} = \left\| y \right\|^{2} + \left\| {f(x)} \right\|^{2} \geq \left\| y \right\|^{2}\), with equality iff \(f(x) = 0\), hence \(\argmin_{x \in F}\left\| {y - f(x)} \right\| = \ker f = \left\{ 0 \right\}\). On the other side \(\max\limits_{\left\| {f(d)} \right\| \leq 1}{\langle y,f(d)\rangle} = 0\), so every maximizer \(d^{\ast}\) satisfies \({\langle y,f\left( d^{\ast} \right)\rangle} = 0\) and the rescaled set \(\left\{ {\langle y,f\left( d^{\ast} \right)\rangle}^{\frac{1}{p}}d^{\ast} \right\} = \left\{ 0 \right\}\). Both sides reduce to \(\left\{ 0 \right\}\), so the proportionality holds.
% \end{itemize}

\end{proof}

\begin{lemma}\label{lemma:optimal_projection}

For:

\begin{itemize}
\item
  \(n,p,a,b,p \in \mathbb{N}\)
\item
  \(Y \in (n,p),X \in (n,a),A \in (a,b),B \in (b,p)\)
\item
  \(s(A):A \mapsto \argmin_{B}\left\| {Y - XAB} \right\|\)
\item
  \(\overline{s}:A \mapsto \overline{s}(A) \in s(A)\)
\item
  \(\hat{Y} \coloneqq XA\overline{s}(A)\)
\end{itemize}

\[\left\| {Y - \hat{Y}} \right\|^{2} = \left\| Y \right\|^{2} - \left\| \hat{Y} \right\|^{2}\]

\end{lemma}

\begin{proof}

By definition of \(\overline{s}(A)\) being optimal, we have that
\(\hat{Y}\) is the orthogonal projection of \(Y\) onto the space spanned
by the columns of \(XA\). Hence \(Y - \hat{Y}\) is orthogonal to
\(\hat{Y}\) and the Pythagorean theorem gives the result.

\end{proof}

\begin{theorem}\label{theorem:separable_optimization}

For:

\begin{itemize}
\item
  \(n,p,a,b,p \in \mathbb{N}\)
\item
  \(Y \in (n,p),X \in (n,a),A \in (a,b),B \in (b,p)\)
\item
  \(s(A):A \mapsto \argmin_{B}\left\| {Y - XAB} \right\|\)
\item
  \(\overline{s}:A \mapsto \overline{s} \in s(A)\)
\end{itemize}

We have that:
\[\argmin_{A,B}\left\| {Y - XAB} \right\| = \bigcup_{A^{\ast}}\left\{ A^{\ast} \right\} \times s\left( A^{\ast} \right)~|~A^{\ast} \in \argmax_{A}\left\| {XA\overline{s}(A)} \right\|\]

\end{theorem}

\begin{proof}

\[
\begin{aligned}
\argmin_{A,B}\left\| {Y - XAB} \right\| & = \argmin_{A}\left\| {Y - XAB} \right\|,B \in s(A)
\end{aligned}
\]
Hence:
\[
\begin{aligned}
    \argmin_{A}\left\| {Y - XAB} \right\| & = \argmin_{A}\left\| {Y - XA\overline{s}(A)} \right\| \\
    \text{By}\ {}: & = \argmin_{A}\left\| Y \right\|^{2} - \left\| \hat{Y} \right\|^{2} \\
     & = \argmax_{A}\left\| {XA\overline{s}(A)} \right\| \\
     & = \bigcup_{A^{\ast}}\left\{ A^{\ast} \right\} \times s(A^{\ast})~|~A^{\ast} \in \argmax_{A}\left\| {XA\overline{s}(A)} \right\|
\end{aligned}
\]

\end{proof}

\begin{theorem}\label{theorem:rimannian_norm}

For:

\begin{itemize}
\item
  \(\left\| . \right\|\) the Frobenius norm
\item
  \(A \in (n,m)\)
\item
  \(X \in (m,n)\) We have that:
  \[\left\| {AX} \right\| = \left\| {\left( A^{\top}A \right)^{\frac{1}{2}}X} \right\|\]
\end{itemize}

\end{theorem}

\begin{proof}

\[\begin{aligned}
\left\| {AX} \right\|^{2} & = {\langle AX,AX\rangle} \\
 & = {\langle A^{\top}AX,X\rangle} \\
 & = {\langle\left( A^{\top}A \right)^{\frac{1}{2}}\left( A^{\top}A \right)^{\frac{1}{2}}X,X\rangle} \\
 & = \left\| {\left( A^{\top}A \right)^{\frac{1}{2}}X} \right\|^{2}
\end{aligned}\]

\end{proof}

\subsubsection{Notations}

\subsubsection{Linear Algebra}

\begin{itemize}
\item
  \(\langle.,.\rangle\) is a scalar product and \(\left\| . \right\|\)
  the associated norm.
\item
  \(\left\| . \right\|_{2}\) is the Frobenius norm.
\item
  \(\times_{k}\) is to indicate multiplication over a common dimension
  \(k\).
\item
  \(.^{+}\) is the Moore-Penrose pseudoinverse.
  \item To make the size of matrices more readable we use $(a, b)$ instead of $\mathbb{R}^{a \times b}$.
\end{itemize}

\subsubsection{Data and Model}

We consider a data distribution \(\D\), a per-data-point loss function
\(\ell:\hat{y} \times y \mapsto \ell(\hat{y},y)\), and a model
\(f:x \times \theta \mapsto f_{\theta}(x) = \hat{y}\) with parameters
\(\theta\) mapping inputs \(x\) to predictions \(\hat{y}\). The overall
loss is:
\[
    \loss(f_{\theta}) = \EE_{(x,y) \sim \D}\left\lbrack \ell(f_{\theta}(x),y) \right\rbrack
\]

\subsubsection{Network Architecture}

We consider a feedforward neural network with layers indexed by \(l\).
For a data point \(x\), the forward pass through consecutive layers is:
\[
    x \overset{a_{l - 2}}{\xrightarrow{\hspace{2em}}} a_{l - 2}(x) \overset{L_{W_{l - 1}}}{\xrightarrow{\hspace{2em}}} s_{l - 1}(x) \overset{\sigma}{\xrightarrow{\hspace{2em}}} a_{l - 1}(x) \overset{L_{W_{l}}}{\xrightarrow{\hspace{2em}}} s_{l}(x) \overset{\hspace{2em}}{\xrightarrow{\hspace{2em}}} f(x)
\]
For readability, we drop layer indices when unambiguous:
\[
    x \overset{a_{- 2}}{\xrightarrow{\hspace{2em}}} a_{- 2}(x) \overset{L_{W_{- 1}}}{\xrightarrow{\hspace{2em}}} s_{- 1}(x) \overset{\sigma}{\xrightarrow{\hspace{2em}}} a_{- 1}(x) \overset{L_{W}}{\xrightarrow{\hspace{2em}}} s(x) \xrightarrow{\hspace{2em}} f(x)
\]

\subsubsection{Dimensions and Batched Notation}

The size of each layer \(l\) is \(C_{l}\) (written as \(C_{- 2}\),
\(C_{- 1}\), \(C\) for layers \(l - 2\), \(l - 1\), \(l\) respectively).
The number of data points is \(n\). For a single sample:
\(a_{- 2}(x) \in \left( C_{- 2},1 \right) = \left( C_{- 2} \right)\). For \(n\) batched samples, we stack row-wise:
\[
    \mathbf{a}_{- 2} = \begin{pmatrix}
    {a_{- 2}(x_{1})}^{\top} \\
    \ldots \\
    {a_{- 2}(x_{n})}^{\top}
    \end{pmatrix} \in \left( n,C_{- 2} \right)
\]
The linear layer acts as: \[\begin{aligned}
\mathrm{L}_{W}(a_{- 2}(x)) & = W_{l - 1} \times_{C_{- 2}}a_{- 2}(x) \in \left( C_{- 1},1 \right) \\
\mathrm{L}_{W}(\mathbf{\mathbf{a}_{- 2}}) & = \mathbf{a}_{- 2} \times_{C_{- 2}}W^{\top} \in \left( n,C_{- 1} \right)
\end{aligned}\]

\subsubsection{Gradient Notation}

The gradient of the loss with respect to pre-activations, stacked
row-wise: \[\Gradpreact = \begin{pmatrix}
{\nabla_{s_{}(x_{1})}\ell(f\left( x_{1} \right))}^{\top} \\
\ldots \\
{\nabla_{s_{}(x_{n})}\ell(f\left( x_{n} \right))}^{\top}
\end{pmatrix} \in (n,C)\]

\subsubsection{Neuron Addition}

We aim to expand the intermediate representation \(s_{- 1}\) and
\(a_{- 1}\) by adding \(\cext\) new neurons. This is done by adding
fan-in weights \(\Fanin{} = \begin{pmatrix}
\fanin{}_{1} \\
\ldots \\
\fanin{}_{k}
\end{pmatrix} \in \left( \cext,C_{- 2} \right)\) and fan-out weights
\(\Fanout{} = \begin{pmatrix}
\fanout{}_{1} & \ldots & \fanout{}_{k}
\end{pmatrix} \in \left( C,\cext \right)\):

\[
a_{- 2}(x) \overset{L_{\Fanin{}}}{\xrightarrow{\hspace{2em}}} \preactext(x) \overset{\sigma}{\xrightarrow{\hspace{2em}}} \postactext(x) \overset{L_{\Fanout{}}}{\xrightarrow{\hspace{2em}}} \delta_s(x)
\]
The change to the pre-activation at layer \(l\) is:
\[s_{} \leftarrow s_{} + \delta_{s_{}}\] with
\(\delta_{s_{}} = \mathrm{L}_{\Fanout{}} \circ \sigma \circ \mathrm{L}_{\Fanin{}} \circ a_{- 2}\).

\subsubsection{Key Statistics}

Several statistics computed over the dataset are used across all
methods:

\begin{definition}

The \textbf{activation covariance} at layer \(l\) is:
\[A_{l} \coloneqq \EE_{(x,y) \sim \D}\left\lbrack a_{l}(x) \times_{1}\left( a_{l}(x) \right)^{\top} \right\rbrack \approx \frac{1}{n}{\mathbf{a}_{l}}^{\top} \times_{n}\mathbf{a}_{l} \;\;\in \left( C_{l},C_{l} \right)\]

\end{definition}

\begin{definition}

The \textbf{gradient-activation cross-covariance} with activation at
layer \(l\) with $l \in \{-1, -2\}$ is:
\[B_{l} \coloneqq \EE_{(x,y) \sim \D}\left\lbrack \left( - {\nabla_{s(x)}\ell(f(x))} \right) \times_{1}\left( a_{l}(x) \right)^{\top} \right\rbrack \approx - \frac{1}{n}\Gradpreact^{\top} \times_{n}\mathbf{a}_{l} \;\;\in \left( C,C_{l} \right)\]

\end{definition}

\begin{definition}

The \textbf{gradient covariance} is:
\[\covgrad \coloneqq \EE_{(x,y) \sim \D}\left\lbrack \left( \frac{\partial f}{\partial s_{}}(x) \right)^{\top} \times_{C_{L}}\frac{\partial f}{\partial s_{}}(x) \right\rbrack\;\; \in (C,C)\]

\end{definition}

\subsubsection{Residual Gradient (Bottleneck)}

To avoid redundancy between added neurons and existing ones, several
methods project the gradient orthogonally to the space spanned by
the activity of existing neurons. The idea is to find the optimal update
\({\mathrm{d}W}^{\ast}\) to the existing weights \(W\) and then consider
only the residual gradient that cannot be addressed by this update.



\begin{definition}

The \textbf{residual gradient} (or \textbf{bottleneck}) is the residual
after subtracting the contribution of the optimal weight update
\({\mathrm{d}W}^{\ast}\):
\[\left( {\nabla_{s_{}}\loss(f)} \right)^{\bot}(x) \coloneqq - {\nabla_{s_{}}\loss(f)}(x) - \mathrm{L}_{{\mathrm{d}W}^{\ast}}(a_{- 1}(x)) \in (C)\]
In batched form:
\[{\Bottleneck} \coloneqq - \Gradpreact - \mathbf{a}_{- 1} \times_{C_{- 1}}A_{+}^{- 1} \times_{C_{- 1}}{B_{- 1}}^{\top} \in (n,C)\]

\end{definition}

\begin{lemma}\label{lemma:optimal_dW}

The optimal weight update \({\mathrm{d}W}^{\ast}\) that minimizes the
residual gradient is:
\[{\mathrm{d}W}^{\ast} \coloneqq \argmin_{\mathrm{d}W}\left\| {- {\nabla_{s_{}}\loss(f)} - \mathrm{L}_{\mathrm{d}W} \circ a_{- 1}} \right\|_{2} = B_{- 1} \times_{C_{- 1}}{A_{- 1}}^{+}\]

\end{lemma}

\begin{proof}

We solve it for a batch of data points:
\({\mathrm{d}W}^{\ast} = \argmin_{\mathrm{d}W}\left\| {- \Gradpreact - \mathbf{a}_{- 1} \times_{C_{- 1}}{\mathrm{d}W}^{\top}} \right\|_{2}\)
This is a least-squares problem:
\[{{\mathrm{d}W}^{\ast}}^{\top} = \left( {\mathbf{a}_{- 1}}^{\top} \times_{n}\mathbf{a}_{- 1} \right)^{+} \times_{C_{- 1}}{\mathbf{a}_{- 1}}^{\top} \times_{n}\left( - \Gradpreact \right) = {A_{- 1}}^{+} \times_{C_{- 1}}{B_{- 1}}^{\top}\]

\end{proof}

\begin{proposition}\label{prop:bottleneck_orthogonality}

The residual gradient \({\Bottleneck}\) is orthogonal to existing
activations:
\[{\mathbf{a}_{- 1}}^{\top} \times_{n}{\Bottleneck} = 0 \in \left( C_{- 1},C \right)\]

\end{proposition}

\begin{proof}

By construction, \(\left( {\nabla_{s_{}}\loss(f)} \right)^{\bot}\) is
the residual of the least-squares projection of
\(- {\nabla_{s_{}}\loss(f)}\) onto the space spanned by \(a_{- 1}\). The
residual is orthogonal to the projection space.

\end{proof}

\subsection{Existing Methods}

This section presents each method with its motivation and derives its
simplified objective function.

\subsubsection{TINY}

\noindent\textbf{Motivation.} The TINY~\citep{verbockhaven_growing_2024} method seeks to approximate the desired pre-activation update \(-\nabla_{s_{}}\loss(f)\) by changes to existing weights and the contribution of new neurons.
It first computes the best update \(\mathrm{d}W^*\) using existing neurons alone, as in \autoref{lemma:optimal_dW}.
Keeping this update fixed, it then chooses new neurons to approximate the remaining residual \(\Bottleneck\).
This is a sequential least-squares construction; it need not solve the joint optimization over existing and new weights~\citep{verbockhaven_growing_2024}.

\noindent\textbf{Objective Formulation.} TINY approximates
\(\delta_{s_{}}(x) \approx \sigma'(0) \cdot \Fanout{} \times_{\cext}\Fanin{} \times_{C_{- 2}}a_{- 2}(x)\)
(linearizing \(\sigma\) around \(0\)). Dropping the constant
\(\sigma'(0)\), the objective becomes:

\begin{proposition}

\textbf{TINY optimization problem (with projection).} The TINY method~\footnote{As stated in  equation (8) of \citep{verbockhaven_growing_2024}.} solves the following problem for the new weights:

\begin{equation}\label{eq:tiny_projection_optimization}
\argmin_{\Fanin{},\Fanout{}}\left\| {{\Bottleneck} - \mathbf{a}_{- 2} \times_{C_{- 2}}\Fanin{}^{\top} \times_{\cext}\Fanout{}^{\top}} \right\|_{2}
\end{equation}
Without projection, replace \({\Bottleneck}\) with \(- \Gradpreact\).

\end{proposition}

\begin{proof}

We maximize the first-order loss decrease. If $a_{-2}^\top G^\bot \neq 0$, the non degeneracy condition is met and by
\autoref{corollary:argmin_norm_to_argmax_scalarp}:
\[\argmin_{\Fanin{},\Fanout{};\left\| \delta_{s_{}} \right\| \leq 1}{\langle\left( {\nabla_{s_{}}\loss(f)} \right)^{\bot},\delta_{s_{}}\rangle} \propto \argmin_{\Fanin{},\Fanout{}}\left\| {- \left( {\nabla_{s_{}}\loss(f)} \right)^{\bot} - \delta_{s_{}}} \right\|\]
Then substituting the linearized \(\delta_{s_{}}\) and using batched
notation gives the result.

If $a_{-2}^\top G^\bot = 0$, in practice $\deltas = 0$ is selected, which is one of the solutions for both problems.


\end{proof}

\subsubsection{GradMax}

\textbf{Motivation.} GradMax~\citep{evci_gradmax_2022} takes a different
approach: initialize new neurons with zero contribution
(\(\delta_{s_{}} = 0\) by setting \(\Fanin{} = 0\)), then maximize the
gradient magnitude that will be available for subsequent training.

\textbf{Objective Formulation.} With \(\Fanin{} = 0\) and \(\sigma(0) = 0\), we have
\(\postactext = 0\), so \(\nabla_{\Fanout{}}\loss(f) = 0\). The loss
decreases after one gradient step on
\(\left( \Fanin{},\Fanout{} \right)\) is:
\[\loss(f_{\Fanin{} + {\mathrm{d}\Fanin{}}, \Fanout{} + {\mathrm{d}\Fanout{}}}) \approx \loss(f) - {\left\| {\nabla_{\Fanin{}}\loss(f)} \right\|_{2}}^{2} - {\left\| {\nabla_{\Fanout{}}\loss(f)} \right\|_{2}}^{2}\]
As \(\nabla_{\Fanout{}}\loss(f) = 0\) at initialization, GradMax
maximizes \(\left\| {\nabla_{\Fanin{}}\loss(f)} \right\|_{2}\).

\begin{proposition}

\textbf{GradMax optimization problem.} GradMax solves~\footnote{As stated in equation (11) of \citep{evci_gradmax_2022}}:
\[\Fanout{}^{\ast} = \argmax_{{\left\| \Fanout{} \right\|_{2}}^{2} \leq \cext}\left\| {\nabla_{\Fanin{}}\loss(f)} \right\|_{2}\]
such that \(\Fanout{}^{\top} \times_{\cext}\Fanout{} = I_{\cext}\).

\end{proposition}

\noindent\textbf{Simplified Objective}

\begin{lemma}

With \(\Fanin{} = 0\) and \(\sigma'(0) = 1\):
\[\nabla_{\Fanin{}}\loss(f) = \Fanout{}^{\top} \times_{C}\EE_{(x,y) \sim \D}\left\lbrack {\nabla_{s_{}}\ell(f(x))} \times_{1}{a_{- 2}(x)}^{\top} \right\rbrack = - \Fanout{}^{\top} \times_{C}B_{- 2}\]

\end{lemma}

\begin{proposition}\label{prop:gradmax_reduced}

\textbf{GradMax reduced objective.} The GradMax optimization reduces to:

\begin{equation}\label{eq:gradmax_reduced}
\mathcal{J}_{\text{GradMax}}\left( \Fanout{} \right) \coloneqq {\left\| {{B_{- 2}}^{\top} \times_{C}\Fanout{}} \right\|_{2}}^{2}
\end{equation}
such that \(\Fanout{}^{\top}\Fanout{} = I_{\cext}\)~\footnote{This constraint is not explicitly mentioned in the paper but is implicit knowing the objective is optimized through SVD.}. The optimal \(\Fanout{}^{\ast}\) are the leading left-singular vectors
of \(B_{- 2}\).

\end{proposition}

\noindent\textbf{GradMax-Opt.} GradMax-Opt is similar, but it does not include the orthogonality constraint. Hence, instead of using an SVD, it uses projected gradient descent. This allows an alternative initialization where the roles of \(\Fanin{}\) and \(\Fanout{}\) are switched: setting \(\Fanout{} = 0\) and maximizing \(\left\| {\nabla_{\Fanout{}}\loss(f)} \right\|_{2}\).

\subsubsection{SENN}\label{sec:senn}

\textbf{Motivation}

SENN~\citep{mitchell_self-expanding_2024} extends GradMax
by considering \textbf{natural gradient} descent instead of standard
gradient descent. It also avoids redundancy by maximizing the
\emph{increase} in natural gradient norm after neuron addition. They
choose to fix \(\Fanout{} = 0\) at initialization to ensure
\(\delta_{s_{}} = 0\) (unlike GradMax that fixes \(\Fanin{} = 0\)).

The natural gradient step is \(\begin{smallpmatrix}
{\mathrm{d}\Fanout{}} \\
{\mathrm{d}W}
\end{smallpmatrix} \propto - F^{- 1}\nabla_{\begin{smallpmatrix}
\Fanout{} \\
W
\end{smallpmatrix}}\loss(f)\) where \(F\) is the Fisher information matrix
with respect to \(\begin{smallpmatrix}
\Fanout{} \\
W
\end{smallpmatrix}\), the first-order equation becomes (also because
\(\nabla_{\Fanin{}}\loss(f) = 0\)):
\[\loss(f_{W + {\mathrm{d}W},\Fanout{} + {\mathrm{d}\Fanout{}}}) \approx \loss(f) - {\langle\nabla_{\begin{smallpmatrix}
\Fanout{} \\
W
\end{smallpmatrix}}\loss(f),F^{- 1}\nabla_{\begin{smallpmatrix}
\Fanout{} \\
W
\end{smallpmatrix}}\loss(f)\rangle}_{2}\] Hence they want to maximize
\(\eta_{\text{SENN}}(\begin{smallpmatrix}
\Fanout{} \\
W
\end{smallpmatrix}) \coloneqq {\langle\nabla_{\begin{smallpmatrix}
\Fanout{} \\
W
\end{smallpmatrix}}\loss(f),F^{- 1}\nabla_{\begin{smallpmatrix}
\Fanout{} \\
W
\end{smallpmatrix}}\loss(f)\rangle}_{2}\).

\noindent\textbf{Objective Formulation}
Using the KFAC~\citep{martens_kfac_2015,dangelKroneckerfactoredApproximateCurvature2025} approximation
\(F \approx \covgrad \times A_{- 1}\) for the Fisher information matrix:

\begin{definition}

The \textbf{SENN criterion} measures the natural gradient norm:
\[\begin{aligned}
\eta_{\text{SENN}}(\begin{smallpmatrix}
\Fanout{} \\
W
\end{smallpmatrix}) & \coloneqq {\langle\nabla_{\begin{smallpmatrix}
\Fanout{} \\
W
\end{smallpmatrix}}\loss(f),F^{- 1}\nabla_{\begin{smallpmatrix}
\Fanout{} \\
W
\end{smallpmatrix}}\loss(f)\rangle}_{2} \\
\text{Using KFAC:}\quad & = {\langle\nabla_{\begin{smallpmatrix}
\Fanout{} \\
W
\end{smallpmatrix}}\loss(f)\left( A_{- 1}^{\text{conc}} \right)^{- 1},\covgrad^{- 1}\nabla_{\begin{smallpmatrix}
\Fanout{} \\
W
\end{smallpmatrix}}\loss(f)\rangle}_{2}
\end{aligned}\] where:

\begin{itemize}
\item
  \(A_{- 1}^{\text{conc}}\) is the covariance of concatenated
  activations \(\begin{smallpmatrix}
  \postactext \\
  a_{- 1}
  \end{smallpmatrix}\)
\end{itemize}

\end{definition}

SENN maximizes the increase
\(\Delta\eta_{\text{SENN}} \coloneqq \eta_{\text{SENN}}(\begin{smallpmatrix}
\Fanout{} \\
W
\end{smallpmatrix}) - \eta_{\text{SENN}}(W)\).

\noindent\textbf{Simplified Objective}
SENN uses a tractable lower bound
\(\Delta\eta_{\text{SENN}}' \leq \Delta\eta_{\text{SENN}}\) involving
the residual gradient.

\begin{definition}
The \textbf{extended activation covariance} (using exact activations):
\[A_{- 1}^{\text{ext}} \coloneqq \EE_{(x,y) \sim \D}\left\lbrack \postactext(x) \times_{1}{\postactext(x)}^{\top} \right\rbrack \in \left( \cext,\cext \right)\]
\end{definition}

\begin{definition}
The \textbf{residual gradient-activation cross-covariance} (using exact
activations):
\[\partial\Fanout{} \coloneqq \EE_{(x,y) \sim \D}\left\lbrack \left( \nabla_{s(x)}\ell(f(x)) \right)^{\bot} \times_{1}{\postactext(x)}^{\top} \right\rbrack \in \left( C,\cext \right)\]
\end{definition}

\begin{proposition}\label{prop:senn_reduced}
\textbf{SENN reduced objective.~\footnote{As stated in Theorem 3.2 of~\citep{mitchell_self-expanding_2024}.}} The SENN lower bound objective is
\(\eta_{\text{SENN}}(\Fanout{})\) but computing the gradient of
\(\Fanout{}\) backpropagating only the residual gradient (hence the
definition of \(\partial\Fanout{}\) the pseudo gradient of
\(\Fanout{}\)):
\begin{equation}\label{eq:senn_reduced}
\begin{aligned}
\mathcal{J}_{\text{SENN}}\left( \Fanin{} \right) & \coloneqq {\langle\partial\Fanout{} \times_{\cext}{A_{- 1}^{\text{ext}}}^{+},\covgrad^{- 1} \times_{C}\partial\Fanout{}\rangle}_{2} \\
 & = \left\| {\isqrtS \times_{C}\partial\Fanout{} \times_{\cext}\isqrtAext} \right\|_{2}^{2}
\end{aligned}
\end{equation}
\end{proposition}

\subsubsection{NeST}

\noindent\textbf{Motivation}

NeST~\citep{daiNeSTNeuralNetwork2019} considers sparse networks where
connections are added incrementally. NeST provides criteria for both
\textbf{connection activation} (adding a connection in an existing
layer) and \textbf{neuron addition} (adding a new neuron with a single
connection pair).

\textbf{Connection Activation}

NeST considers sparse layers where \(W\) is sparse. To choose which
connection to add, they look at which connection between neuron \(i\) in
layer \(l - 1\) and neuron \(j\) in layer \(l\) would maximally decrease
the loss if added and optimized with a gradient step:
\[\loss(f_{W + {\mathrm{d}W}\lbrack i,j\rbrack}) = \loss(f) + {\langle\nabla_{W\lbrack i,j\rbrack}\loss(f),{\mathrm{d}W}\lbrack i,j\rbrack\rangle} + o\left( \left\| {{\mathrm{d}W}\lbrack i,j\rbrack} \right\|_{2} \right)\]
When choosing
\({\mathrm{d}W}\lbrack i,j\rbrack \propto - \nabla_{W\lbrack i,j\rbrack}\loss(f)\),
this becomes:
\[\loss(f_{W + {\mathrm{d}W}\lbrack i,j\rbrack}) = \loss(f) - {\left\| {\nabla_{W\lbrack i,j\rbrack}\loss(f)} \right\|_{2}}^{2} + o\left( \left\| {{\mathrm{d}W}\lbrack i,j\rbrack} \right\|_{2} \right)\]
Note that NeST justifies this choice with Hebbian theory and considers
\(\left| {\nabla_{W\lbrack i,j\rbrack}\loss(f)} \right|\).

\begin{proposition}\label{prop:nest_connection}

\textbf{NeST connection activation.} For adding a connection in an
existing layer, NeST selects:
\[i^{\ast},j^{\ast} = \argmax_{i,j}\left| B_{- 1} \right|\lbrack j,i\rbrack\]

\end{proposition}

\noindent\textbf{Neuron Addition}
When adding a new neuron, NeST only adds a few connections between pairs
of neurons \((i,j)\) with neuron \(i\) in layer \(l - 2\) and neuron
\(j\) in layer \(l\). They set \(\Fanin{}\) and \(\Fanout{}\) to zero
except for the values at indices \((i)\) and \((j)\) respectively. This
gives:
\[\delta_{s_{}}(x)\lbrack j\rbrack = \Fanout{}\lbrack j\rbrack\sigma(\Fanin{}\lbrack i\rbrack a_{- 2}(x)\lbrack i\rbrack) \approx \Fanout{}\lbrack j\rbrack\sigma'(0)\Fanin{}\lbrack i\rbrack a_{- 2}(x)\lbrack i\rbrack\]
Dropping the constant \(\sigma'(0)\), they optimize the decrease in loss
after the addition: \[\begin{aligned}
\loss(s_{} + \delta_{s_{}}) & = \loss(s_{}) + {\langle\nabla_{s_{}\lbrack j\rbrack}\loss(s_{}),\delta_{s_{}}\lbrack j\rbrack\rangle} + o\left( \left\| {\delta_{s_{}}\lbrack j\rbrack} \right\| \right) \\
 & \approx \loss(s_{}) + \nabla_{s_{}\lbrack j\rbrack}\loss(s_{}) \cdot \Fanout{}\lbrack j\rbrack \cdot \Fanin{}\lbrack i\rbrack \cdot a_{- 2}\lbrack i\rbrack
\end{aligned}\]
Hence they minimize
\(\nabla_{s_{}\lbrack j\rbrack}\loss(s_{}) \cdot a_{- 2}\lbrack i\rbrack \cdot \Fanout{}\lbrack j\rbrack \cdot \Fanin{}\lbrack i\rbrack\)
under the constraint that
\(\left| {\Fanout{}\lbrack j\rbrack \cdot \Fanin{}\lbrack i\rbrack} \right|\)
is small.

\begin{proposition}\label{prop:nest_objective}

\textbf{NeST neuron addition.} For adding a new neuron, NeST selects:
\begin{equation}\label{eq:nest_neuron_metric}
i^{\ast},j^{\ast} = \argmax_{i,j}\left| B_{- 2} \right|\lbrack j,i\rbrack
\end{equation}
to add a single connection, or the pairs in the top \(\beta\) fraction
to add several connections.

\end{proposition}

\noindent\textbf{Initialization} Once \(\left( i^{\ast},j^{\ast} \right)\) is chosen, NeST initializes the weights so that the change \(\delta_{s_{}}\lbrack j\rbrack\) is similar to that created by a gradient step on a virtual connection
between \(a_{- 2}\lbrack i\rbrack\) and \(s_{}\lbrack j\rbrack\), yielding a loss decrease of \({\left\| {B_{- 2}\left\lbrack j^{\ast},i^{\ast} \right\rbrack} \right\|_{2}}^{2}\). The sign is randomized:

\begin{proposition}\label{prop:nest_init}

\textbf{NeST initialization.} Once \(\left( i^{\ast},j^{\ast} \right)\)
is selected, NeST initializes: \[\begin{aligned}
\varepsilon & \sim \text{ Rademacher }\text{i.e.}\text{Uniform}\left( \left\{ - 1,1 \right\} \right) \\
\Fanin{}\left\lbrack i^{\ast} \right\rbrack & = \varepsilon \cdot \text{ sign}\left( B_{- 2}\left\lbrack j^{\ast},i^{\ast} \right\rbrack \right) \cdot \sqrt{\left| {B_{- 2}\left\lbrack j^{\ast},i^{\ast} \right\rbrack} \right|} \\
\Fanout{}\left\lbrack j^{\ast} \right\rbrack & = \varepsilon \cdot \sqrt{\left| {B_{- 2}\left\lbrack j^{\ast},i^{\ast} \right\rbrack} \right|}
\end{aligned}\]
All other entries of \(\Fanin{}\) and \(\Fanout{}\) are zero.

\textbf{Note:} The NeST paper seems to align the product
\(\Fanout{}\left\lbrack j^{\ast} \right\rbrack \cdot \Fanin{}\left\lbrack i^{\ast} \right\rbrack\)
with \(- B_{- 2}\left\lbrack j^{\ast},i^{\ast} \right\rbrack\) instead
of \(B_{- 2}\left\lbrack j^{\ast},i^{\ast} \right\rbrack\), which would
maximize the loss instead of minimizing it. We present here what we
believe is the intended version.

\end{proposition}

\subsubsection{NORTH}

\textbf{Motivation}

NORTH~\citep{maile_when_2022} takes a geometric perspective:
new neurons should add \textbf{new directions} in activation space,
orthogonal to existing activations.

\textbf{Objective Formulation}

NORTH seeks \(\fanin{}\) such that the new activations
\(\Postactext = \sigma(\mathbf{a}_{- 2}\fanin{}^{\top})\) are orthogonal
to existing ones, which is equivalent to lying in
\(\ker({\mathbf{a}_{- 1}}^{\top})\).

\begin{proposition}\label{prop:north_construction}

\textbf{NORTH-Pre construction.} Let
\(V \in \left( n,n - C_{- 1} \right)\) be an orthonormal basis of
\(\ker({\mathbf{s}_{- 1}}^{\top})\) (computed via SVD). For random
directions \(R \in \left( n - C_{- 1},\cext \right)\):

\begin{equation}\label{eq:north_proposition}
\Fanin{}^{\top} = A_{- 2}^{- 1} \times_{C_{- 2}}{\mathbf{a}_{- 2}}^{\top} \times_{n}V \times_{n - C_{- 1}}R
\end{equation}

\end{proposition}

\paragraph{Justification}

They use the approximation \(\sigma \approx \text{ Identity}\) to relax
the problem to replace \(\Postactext\) with
\(\Preactext = \mathbf{a}_{- 2} \times_{C_{- 2}}\fanin{}^{\top}\) and
\(\ker({\mathbf{a}_{- 1}}^{\top})\) with
\(\ker({\mathbf{s}_{- 1}}^{\top})\).

First they compute a basis \(V \in \left( n,n - C_{- 1} \right)\) of
\(\ker({\mathbf{s}_{- 1}}^{\top})\) using SVD. Then they generate random
combinations of the basis vectors \(V \times_{n - C_{- 1}}r\) for random
\(r\) as candidates for \(\Preactext\). For each candidate, they solve
for \(\fanin{}\) such that
\(\mathbf{a}_{- 2} \times_{C_{- 2}}\fanin{}^{\top} = V \times_{n - C_{- 1}}r\)
which gives:
\[\fanin{}^{\top} = \left( \mathbf{a}_{- 2} \right)^{+} \times_{n}Vr\]

Or for a batched version with \(\cext\) candidates
\(R \in \left( n - C_{- 1},\cext \right)\) with orthogonal columns:
\[\begin{aligned}
\Fanin{}^{\top} & = \left( \mathbf{a}_{- 2} \right)^{+} \times_{n}V \times_{n - C_{- 1}}R \in \left( C_{- 2},\cext \right) \\
 & = \left( {\mathbf{a}_{- 2}}^{\top} \times_{n}\mathbf{a}_{- 2} \right)^{+} \times_{C_{- 2}}{\mathbf{a}_{- 2}}^{\top} \times_{n}V \times_{n - C_{- 1}}R \\
 & \propto A_{- 2}^{- 1} \times_{C_{- 2}}{\mathbf{a}_{- 2}}^{\top} \times_{n}V \times_{n - C_{- 1}}R
\end{aligned}\] (\(\propto\) because we skip the \(\frac{1}{n}\)
factor.) Then they select the best candidates according to the
orthogonality of
\(\Postactext = \sigma(\mathbf{a}_{- 2}\fanin{}^{\top})\) to
\(\mathbf{a}_{- 1}\) using this time the exact activation function
\(\sigma\).

\subsection{Unified Framework}

For completeness, we now restate \autoref{thm:unified_obj_main_text}, which shows that all methods can be understood as special cases of a
common optimization objective.

\begin{theorem}\label{thm:unified_objective}

\textbf{Unified objective for neuron addition.} All neuron addition
methods studied in this document can be derived from the following
optimization problem:

\begin{subequations}\label{eq:unified_objective}
\begin{align}
\Fanin{}^{\ast},\Fanout{}^{\ast} & = \argmax_{\Fanin{},\Fanout{},\left\| {\sqrtAext \times_{\cext}\Fanout{}^{\top} \times_{C}\sqrtS} \right\|_{2} \leq 1}{\langle T \times_{C}\isqrtS,\frac{1}{n}\Postactext \times_{\cext}\Fanout{}^{\top} \times_{C}\sqrtS\rangle} \label{eq:unified_objective_scalarp} \\
\shortintertext{if we linearize $\sigma$ around $0$ and $\transp{\postact[-2]}T \neq 0$:}
 & \propto \argmin_{\Fanin{},\Fanout{}}\left\| {T \times_{C}\isqrtS - \sigma(\mathbf{a}_{- 2} \times_{C_{- 2}}\Fanin{}^{\top}) \times_{\cext}\Fanout{}^{\top} \times_{C}\sqrtS} \right\|_{2} \label{eq:unified_objective_norm}
\end{align}
\end{subequations}
where:
\begin{itemize}
\item
  \(T \in (n,C)\) is the target (negative gradient \(- \Gradpreact\) or
  residual gradient \({\Bottleneck}\))
\item
  \(\covgrad \in (C,C)\) is the gradient weighting metric
\item
  \(\sigma\) is an activation function
\item
  Additional constraints may be imposed on \(\Fanin{}\) and
  \(\Fanout{}\)
\end{itemize}

Note that if $\transp{\postact[-2]}T = 0$ (under the linearization assumption), $(0, 0)$ is an optimal solution in both cases and is the one selected in practice.
The following table summarizes the specialization for each method:
\begin{longtable}[]{@{}ccccccccc@{}}
\caption{Neuron addition methods as special cases of the unified
objective in \autoref{eq:unified_objective}.}\tabularnewline
\toprule\noalign{}
\textbf{Method} & \(T\) & \(A_{- 1}^{\text{ext}}\) & \(\covgrad\) &
\textbf{Extra constraints} & \(\Fanin{}\) & \(\Fanout{}\) &
\(\postactext \perp a_{- 1}\) & \(\sigma\) \\
\midrule\noalign{}
\endfirsthead
\toprule\noalign{}
\textbf{Method} & \(T\) & \(A_{- 1}^{\text{ext}}\) & \(\covgrad\) &
\textbf{Extra constraints} & \(\Fanin{}\) & \(\Fanout{}\) &
\(\postactext \perp a_{- 1}\) & \(\sigma\) \\
\midrule\noalign{}
\endhead
\bottomrule\noalign{}
\endlastfoot
Full & \({\Bottleneck}\) & \(A_{- 1}^{\text{ext}}\) & \(\covgrad\) &
\(\varnothing\) & \(\Fanin{}^{\ast}\) & \(\Fanout{}^{\ast}\) &
\(\checkmark\) & Any \\
TINY (proj.) & \({\Bottleneck}\) & \(\Fanin{}A_{- 2}\Fanin{}^{\top}\) &
\(I_{C}\) & \(\varnothing\) & \(\Fanin{}^{\ast}\) & \(\Fanout{}^{\ast}\)
& \(\checkmark\) & Identity \\
TINY (no proj.) & \(- \Gradpreact\) & \(\Fanin{}A_{- 2}\Fanin{}^{\top}\)
& \(I_{C}\) & \(\varnothing\) & \(\Fanin{}^{\ast}\) &
\(\Fanout{}^{\ast}\) &
\(\times\) & Identity \\
GradMax & \(- \Gradpreact\) & \(\Fanin{}\Fanin{}^{\top}\) & \(I_{C}\) &
\(\Fanin{}\Fanin{}^{\top} = I_{\cext};\Fanout{}^{\top}\Fanout{} = I_{\cext}\)
& \(0\) & \(\Fanout{}^{\ast}\) &
\(\times\) & Any \\
SENN & \({\Bottleneck}\) & \(A_{- 1}^{\text{ext}}\) & \(\covgrad\) &
\(\varnothing\) & \(\Fanin{}^{\ast}\) & \(0\) &
\(\checkmark\) & Any \\
NeST & \(- \Gradpreact\) & \(\Fanin{}\Fanin{}^{\top}\) & \(I_{C}\) &
\(\left\| \Fanin{} \right\|_{0} = \left\| \Fanout{} \right\|_{0} = 1\) &
\(\Fanin{}^{\ast}\) & \(\Fanout{}^{\ast}\) &
\(\times\) & Identity \\
\end{longtable}

\phantomsection\label{tab:comparison}

\noindent Notes:

\begin{itemize}
\item
  Assuming \(A_{- 1}^{\text{ext }} = \Fanin{}A_{- 2}\Fanin{}^{\top}\) in
  \autoref{eq:unified_objective_scalarp} is equivalent to linearizing
  \(\sigma\) around \(0\) in \autoref{eq:unified_objective_norm}.
\item
  The GradMax assumption
  \(A_{- 1}^{\text{ext }} = \Fanin{}\Fanin{}^{\top}\) is like assuming
  the input data is whitened in addition to linearizing \(\sigma\)
  around \(0\).
\item
  The equivalence between the two optimization problems is correct
  either if we use the first-order approximation of \(\sigma\) or if we
  allow to have different proportionality constants for \(\Fanin{}\) and
  \(\Fanout{}\).
\end{itemize}

\end{theorem}

\begin{proof}

The derivations for each method are provided in the following sections.
We only prove the equivalence between the minimization and maximization
forms. By \autoref{lemma:argmin_norm_to_argmax_scalarp} from the
preliminaries, we have:
\begin{equation}
\begin{aligned}
\Fanin{}^{\ast},\Fanout{}^{\ast} & = \argmin_{\Fanin{},\Fanout{}}\left\| {T \times_{C}\isqrtS - \sigma(\mathbf{a}_{- 2} \times_{C_{- 2}}\Fanin{}^{\top}) \times_{\cext}\Fanout{}^{\top} \times_{C}\sqrtS} \right\|_{2} \\
 & = \argmin_{\Fanin{},\Fanout{}}\frac{1}{\sqrt{n}}\left\| {T \times_{C}\isqrtS - \Postactext \times_{\cext}\Fanout{}^{\top} \times_{C}\sqrtS} \right\|_{2}
\end{aligned}
\end{equation}
Linearizing $\sigma$ around $0$, $\Postactext = \mathbf{a}_{- 2} \times_{C_{- 2}}\Fanin{}^{\top}$ is linear in $\Fanin{}$, so $f:\left( \Fanin{},\Fanout{} \right) \mapsto \frac{1}{\sqrt{n}}\Postactext\Fanout{}^{\top}\sqrtS$ is homogeneous of degree $2$, $n:\left( \Fanin{},\Fanout{} \right) \mapsto \left( \Fanin{}, - \Fanout{} \right)$ is a valid sign-flip operator. As we assume $\transp{\postact[-2]}T \neq 0$ the non-degeneracy condition is met and \autoref{corollary:argmin_norm_to_argmax_scalarp} applies:
\begin{equation}
\begin{aligned}
 & \propto \argmax_{\Fanin{},\Fanout{},\left\| {\frac{1}{\sqrt{n}}\Postactext \times_{\cext}\Fanout{}^{\top} \times_{C}\sqrtS} \right\|_{2} \leq 1}{\langle\frac{1}{\sqrt{n}}T \times_{C}\isqrtS,\frac{1}{\sqrt{n}}\Postactext \times_{\cext}\Fanout{}^{\top} \times_{C}\sqrtS\rangle} \\
 & = \argmax_{\Fanin{},\Fanout{},\left\| {\sqrtAext \times_{\cext}\Fanout{}^{\top} \times_{C}\sqrtS} \right\|_{2} \leq 1}{\langle T \times_{C}\isqrtS,\frac{1}{n}\Postactext \times_{\cext}\Fanout{}^{\top} \times_{C}\sqrtS\rangle}
\end{aligned}
\end{equation}

the last step using \autoref{theorem:rimannian_norm} together with
\(A_{- 1}^{\text{ext }} = \frac{1}{n}\Postactext^{\top} \times_{n}\Postactext\).
This equivalence is only needed for the norm-form methods, i.e., TINY,
which assumes \(\sigma \approx \text{ Identity}\); GradMax, NeST and
SENN are obtained directly from the maximization form
\autoref{eq:unified_objective_scalarp}.

\end{proof}

\subsubsection{Derivation of Each Method}

We now prove that each method can be derived from the unified objective
\autoref{eq:unified_objective} by appropriate specialization.

\noindent\textbf{Derivation of TINY}

\begin{proposition}\label{prop:tiny_derivation}

\textbf{TINY from unified objective.} Setting \(\covgrad = I_{C}\) and
\(\sigma \approx \text{ Identity}\) in \autoref{eq:unified_objective_norm}
yields the TINY optimization problem.

\end{proposition}

\begin{proof}

With \(\covgrad = I_{C}\), we have \(\isqrtS = I_{C}\). The unified
objective becomes:
\[\argmin_{\Fanin{},\Fanout{}}\left\| {T - \sigma(\mathbf{a}_{- 2} \times_{C_{- 2}}\Fanin{}^{\top}) \times_{\cext}\Fanout{}^{\top}} \right\|_{2}\]
Using the linear approximation \(\sigma \approx \text{ Identity}\)
(which is equivalent to
\(A_{- 1}^{\text{ext }} = \Fanin{}A_{- 2}\Fanin{}^{\top}\)), this
becomes:
\[\argmin_{\Fanin{},\Fanout{}}\left\| {T - \mathbf{a}_{- 2} \times_{C_{- 2}}\Fanin{}^{\top} \times_{\cext}\Fanout{}^{\top}} \right\|_{2}\]
With \(T = {\Bottleneck}\), this is exactly
\autoref{eq:tiny_projection_optimization}. With \(T = - \Gradpreact\), this
is TINY without projection.

\end{proof}

\noindent\textbf{Derivation of GradMax}

\begin{proposition}\label{prop:gradmax_derivation}

\textbf{GradMax from unified objective.} Setting
\(T \leftarrow - \Gradpreact\), \(A_{- 2} \leftarrow I_{C_{- 2}}\),
\(\covgrad \leftarrow I_{C}\), with the constraint
\(\Fanout{}^{\top}\Fanout{} = I_{\cext}\), and then setting
\(\Fanin{} = 0\) after optimization yields GradMax.
\end{proposition}
This proof is largely inspired by the work of~\citep{verbockhaven_growing_2024}.

\begin{proof}

Starting from \autoref{eq:unified_objective_scalarp} with
\(\covgrad \leftarrow I_{C}\) and \(T \leftarrow - \Gradpreact\), the
objective is: \[\begin{aligned}
\argmax_{\Fanin{},\Fanout{},\left\| {\left( \Fanin{}\Fanin{}^{\top} \right)^{\frac{1}{2}}\Fanout{}^{\top}} \right\|_{2} \leq 1} & {\langle - \Gradpreact,\Postactext \times_{\cext}\Fanout{}^{\top}\rangle}
\end{aligned}\] Using \autoref{theorem:rimannian_norm} with
\(A \leftarrow \Fanin{}^{\top}\) we have that
\(\left\| {\left( \Fanin{}\Fanin{}^{\top} \right)^{\frac{1}{2}}\Fanout{}^{\top}} \right\|_{2} = \left\| {\Fanin{}^{\top}\Fanout{}^{\top}} \right\|_{2}\).
In addition, we replace \(\Postactext\) with its linear approximation
\(\mathbf{a}_{- 2}\Fanin{}^{\top}\). However, this approximation is not a
strong assumption, as it is only used during backward at
\(\sigma(\mathbf{a}_{- 2}\Fanin{}^{\top}) = \sigma(0)\) as GradMax set
\(\Fanin{} = 0\). Therefore as long as \(\sigma'(0) = 1\) the
approximation is exact. We get:

\[\begin{aligned}
 = \argmax_{\left\| {\Fanin{}^{\top} \times_{\cext}\Fanout{}^{\top}} \right\|_{2} \leq 1} & {\langle - \Gradpreact,\frac{1}{n}\mathbf{a}_{- 2}\Fanin{}^{\top} \times_{\cext}\Fanout{}^{\top}\rangle} \\
 = \argmax_{\left\| {\Fanin{}^{\top} \times_{\cext}\Fanout{}^{\top}} \right\|_{2} \leq 1} & {\langle{\frac{1}{n}\mathbf{a}_{- 2}}^{\top} \times_{n}\left( - \Gradpreact \right),\Fanin{}^{\top} \times_{\cext}\Fanout{}^{\top}\rangle} \\
 = \argmax_{\left\| {\Fanin{}^{\top} \times_{\cext}\Fanout{}^{\top}} \right\|_{2} \leq 1} & {\langle{B_{- 2}}^{\top},\Fanin{}^{\top} \times_{\cext}\Fanout{}^{\top}\rangle} \\
\text{Using}\ {}: \propto \argmin_{\Fanin{},\Fanout{}} & \left\| {{B_{- 2}}^{\top} - \Fanin{}^{\top} \times_{\cext}\Fanout{}^{\top}} \right\|_{2}
\end{aligned}\]
This is equivalent to a low-rank approximation of \({B_{- 2}}^{\top}\).
By \autoref{theorem:separable_optimization}, optimizing over \(\Fanout{}\)
with
\(\Fanin{}^{\top} \leftarrow {\Fanin{}^{\ast}}^{\top} = {B_{- 2}}^{\top}\Fanout{}\left( \Fanout{}^{\top}\Fanout{} \right)^{+} = {B_{- 2}}^{\top}\Fanout{}\)
gives: \[\begin{aligned}
 & \propto \argmax_{{\left\| \Fanout{} \right\|_{2}}^{2} \leq \cext}\left\| {{B_{- 2}}^{\top}\Fanout{}\Fanout{}^{\top}} \right\|_{2} \\
\text{Using}\ {\Fanout{}^{\top}\Fanout{} = I_{\cext}}:\qquad & = \argmax_{{\left\| \Fanout{} \right\|_{2}}^{2} \leq \cext}\left\| {{B_{- 2}}^{\top}\Fanout{}} \right\|_{2}
\end{aligned}\]
This is exactly the GradMax objective
\(\mathcal{J}_{\text{GradMax}}\left( \Fanout{} \right)\) from
\autoref{eq:gradmax_reduced}. Finally, setting \(\Fanin{} = 0\) after
optimization matches GradMax's initialization.

\end{proof}

\noindent\textbf{Derivation of SENN}

\begin{proposition}\label{prop:senn_derivation}

\textbf{SENN from unified objective.} The SENN objective
\(\mathcal{J}_{\text{SENN}}\) is recovered from the maximization form
\autoref{eq:unified_objective_scalarp} with \(T = {\Bottleneck}\) and the
full gradient covariance \(\covgrad\), setting \(\Fanout{} = 0\) after
optimization.

\end{proposition}

\begin{proof}

We start from the maximization form \autoref{eq:unified_objective_scalarp}
with \(T = {\Bottleneck}\) and the full metric \(\covgrad\). Since
\({\langle{\Bottleneck} \times_{C}\isqrtS,X \times_{C}\sqrtS\rangle} = {\langle{\Bottleneck},X\rangle}\)
for any \(X\), the \(\covgrad\) factors cancel in the objective (the
metric only enters through the constraint), which becomes
\[\argmax_{\Fanin{},\Fanout{},\left\| {\sqrtAext \times_{\cext}\Fanout{}^{\top} \times_{C}\sqrtS} \right\|_{2} \leq 1}\frac{1}{n}{\langle{\Bottleneck},\Postactext \times_{\cext}\Fanout{}^{\top}\rangle}.\]
Fix \(\Fanin{}\) and let
\(M \coloneqq \sqrtAext \times_{\cext}\Fanout{}^{\top} \times_{C}\sqrtS\),
so the constraint reads \(\left\| M \right\|_{2} \leq 1\) and
\(\Fanout{}^{\top} = \isqrtAext \times_{\cext}M \times_{C}\isqrtS\).
Moving \(\Postactext\) across the inner product brings out the
cross-covariance
\(\partial\Fanout{} = \frac{1}{n}{\Bottleneck}^{\top} \times_{n}\Postactext\),
and since \(\isqrtAext\) and \(\isqrtS\) are symmetric,
\[\begin{aligned}
\frac{1}{n}{\langle{\Bottleneck},\Postactext \times_{\cext}\Fanout{}^{\top}\rangle} & = {\langle\frac{1}{n}\Postactext^{\top} \times_{n}{\Bottleneck},\Fanout{}^{\top}\rangle} = {\langle{\partial\Fanout{}}^{\top},\Fanout{}^{\top}\rangle} \\
 & = {\langle{\partial\Fanout{}}^{\top},\isqrtAext \times_{\cext}M \times_{C}\isqrtS\rangle} \\
 & = {\langle\isqrtAext \times_{\cext}{\partial\Fanout{}}^{\top} \times_{C}\isqrtS,M\rangle},
\end{aligned}\] whose maximum over \(\left\| M \right\|_{2} \leq 1\) is,
by Cauchy--Schwarz (attained at
\(M \propto \isqrtAext{\partial\Fanout{}}^{\top}\isqrtS\)),
\[\left\| {\isqrtAext \times_{\cext}{\partial\Fanout{}}^{\top} \times_{C}\isqrtS} \right\|_{2} = \left\| {\isqrtS \times_{C}\partial\Fanout{} \times_{\cext}\isqrtAext} \right\|_{2} = \sqrt{\mathcal{J}_{\text{SENN}}\left( \Fanin{} \right)}\]
with \(\mathcal{J}_{\text{SENN}}\) as in \autoref{eq:senn_reduced}. The
outer maximization over \(\Fanin{}\) is thus
\(\argmax_{\Fanin{}}\mathcal{J}_{\text{SENN}}\left( \Fanin{} \right)\),
the SENN objective. As in \autoref{prop:senn_reduced} the optimal
\(\Fanout{}\) only serves to score \(\Fanin{}\); SENN keeps
\(\Fanout{} = 0\) at initialization to preserve the function.

\end{proof}

\noindent\textbf{Derivation of NeST}

\begin{proposition}\label{prop:nest_derivation}

\textbf{NeST from unified objective.} Setting
\(T \leftarrow - \Gradpreact\), \(\covgrad \leftarrow I_{C}\),
linearizing \(\sigma\) and adding the sparsity constraint
\(\left\| \Fanin{} \right\|_{0} = \left\| \Fanout{} \right\|_{0} = 1\)
yields NeST.

\end{proposition}

\begin{proof}

Starting from \autoref{eq:unified_objective_scalarp} with
\(A_{- 1}^{\text{ext }} \leftarrow \Fanin{}\Fanin{}^{\top}\),
\(\covgrad \leftarrow I_{C}\), \(T \leftarrow - \Gradpreact\) and
\(\sigma = \text{Id}\), the objective is: \[\begin{aligned}
\argmax_{\Fanin{},\Fanout{},\left\| {\left( \Fanin{}\Fanin{}^{\top} \right)^{\frac{1}{2}}\Fanout{}^{\top}} \right\|_{2} \leq 1,\left\| \Fanin{} \right\|_{0} = \left\| \Fanout{} \right\|_{0} = 1} & {\langle - \Gradpreact,\mathbf{a}_{- 2} \times_{C_{- 2}}\Fanin{}^{\top} \times_{\cext}\Fanout{}^{\top}\rangle} \\
 = \argmax_{\Fanin{},\Fanout{},\left\| {\Fanin{}^{\top}\Fanout{}^{\top}} \right\|_{2} \leq 1,\left\| \Fanin{} \right\|_{0} = \left\| \Fanout{} \right\|_{0} = 1} & {\langle{B_{- 2}}^{\top},\Fanin{}^{\top} \times_{\cext}\Fanout{}^{\top}\rangle}
\end{aligned}\]
To maximize the decrease in the objective, we choose:
\[i^{\ast},j^{\ast} = \argmax_{i,j}\left| {B_{- 2}\lbrack j,i\rbrack} \right|\]
and set:
\(\Fanin{}\left\lbrack i^{\ast} \right\rbrack = \varepsilon{\operatorname{sign}(B_{- 2}\left\lbrack j^{\ast},i^{\ast} \right\rbrack)}\sqrt{\left| {B_{- 2}\left\lbrack j^{\ast},i^{\ast} \right\rbrack} \right|}\)
and
\(\Fanout{}\left\lbrack j^{\ast} \right\rbrack = \varepsilon\sqrt{\left| {B_{- 2}\left\lbrack j^{\ast},i^{\ast} \right\rbrack} \right|}\)
for random \(\varepsilon \in \left\{ - 1,1 \right\}\). This is exactly
the NeST selection criterion from \autoref{prop:nest_objective}. (Note that
we don't respect the exact scaling constraint
\(\left\| {\Fanin{}^{\top}\Fanout{}^{\top}} \right\|_{2} \leq 1\), but
everything is done up to scaling, so this is not an issue.)

\end{proof}

\subsubsection{NORTH Satisfaction}

\begin{proposition}\label{prop:north_satisfaction}

\textbf{Implicit orthogonality in projected methods.} Methods using the
residual gradient \(T = {\Bottleneck}\) implicitly encourage new
activations to be orthogonal to existing ones, similar to NORTH's
explicit construction.

\end{proposition}

\begin{proposition}

\(\fanin{}\) selected by any method with \(T \leftarrow {\Bottleneck}\)
and \(A_{- 1}^{\text{ext }} \leftarrow \Fanin{}A_{- 2}\Fanin{}^{\top}\)
are in the set of \(\fanin{}\) constructed by NORTH-Pre.

\end{proposition}

\begin{proof}

By \autoref{prop:bottleneck_orthogonality},
\({\mathbf{a}_{- 1}}^{\top} \times_{n}{\Bottleneck} = 0\), so
\({\Bottleneck}\) spans directions orthogonal to \(\mathbf{a}_{- 1}\),
similarly to \(V\), which spans directions orthogonal to
\(\mathbf{s}_{- 1}\) (the same under the first-order approximation,
though further from NORTH's actual goal).

For those methods, the optimal solution satisfies:
\[\mathbf{a}_{- 2} \times_{C_{- 2}}\Fanin{}^{\top} \times_{\cext}\Fanout{}^{\top} \approx {\Bottleneck}\]
Hence (\(\propto\) because we skip the \(\frac{1}{n}\) factor):
\[\Fanin{}^{\top} \propto {A_{- 2}}^{+} \times_{C_{- 2}}{\mathbf{a}_{- 2}}^{\top} \times_{n}{\Bottleneck} \times_{C}{\Fanout{}^{\top}}^{+}\]
Comparing with NORTH-Pre's construction in \autoref{eq:north_proposition}:
\[\Fanin{}^{\top} = A_{- 2}^{- 1} \times_{C_{- 2}}{\mathbf{a}_{- 2}}^{\top} \times_{n}V \times_{n - C_{- 1}}R\]
The structures are identical with \(V \leftarrow {\Bottleneck}\) and
\(R \leftarrow {\Fanout{}^{\top}}^{+}\).

\end{proof}


%%%%%%%%%%%%%%%%%%%%%%%%%%%%%%%%%%%%%%%%%%%%%%%%%%%%%%%%%%%%%%%%%%%%%%%%%%%%%%%%%%%%%%%%%%%%%%%%%%%%%%%%%%%%%%%%%%%%%%%%%%%%%%%%%%%%%%%%%%%%%%%%%%%%%%%%%%%%%%%%%%%%%%%%%%%%%%%%%%%%%%%%%%%%%%%%%%%%%%%%%%%%%%%%%%%%%%%%%%%%%%%%%%%%%%%%%%%%%%%%%%%%%%%%%%%%%%%%%%%%%%%%%%%%%%%%%%%%%%%%%%%%%%%%%%%%%%%%%%%%%%%%%%%%%%%%%%%%%%%%%%%%%%%%%%%%%%%%%%%%%%%%%%%%%%%%%%%%%%%%%%%%%%%%%%%%%%%%%%%%%%%%%%%%%%
